% !TEX root = main.tex
\section{공개 DPP 검증}
DPP 검증은 외부 문서의 지정 필드와 수락된 발급의 연결을 확
인한다. 입력은 알려진 블록체인·원장 문맥, 검증 기준의 확정
시점 \(H_{v}\), DPP 문서와 그 안의 DPPClaim \(=(p,h,c)\)이다. 검
증자는 다음 순서로 판정한다.

1. 배포 문맥과 \(H_{v}\)의 유효성을 확인하고, 지정 규칙에 따라 제

품명(ProductName), 로트 식별자(LotID), 단위(Unit)

를 추출한다. 필드 및 제품 주장 튜플의 정규형을 검사한

다.

2. 문서 해시 \(d\)와 \(H_{Issue}(d,p,c)\)를 재계산하여 제시된 참조

값 \(h\)와 같은지 확인한다.

3. \(H_{v}\)에서 \(h\)가 등록되어 있는지 확인하고, ClaimRef\((h)\)의

생성 기록 색인으로 감사 기록을 읽는다.

4. 생성 기록의 이벤트가 발급이고 정책이 \(p\)이며, 출력 목록

이 정확히 \([\)ClaimRef\((h)]\)인지 확인한다.

어느 조건이든 실패하면 검증을 거부한다. 정책의 공개 의미
를 함께 확인해야 해당 제품 주장이 어떤 기준을 통과했다는 것
인지 해석할 수 있다. 정상적인 수락 관계를 전제로 유효한 제
품 주장은 발급 당시 정책을 만족한 회계 증표가 지정 문서 필드
에 결합되었다는 의미다. 제품 주장에는 상태가 없으므로 현재
활성 여부를 조회하지 않는다. 정책의 사후 비활성화도 과거 제
품 주장을 무효화하지 않는다.

공개 논스 \(c\)는 접근 권한 증표가 아니다. 발급 증명 안에서 논
스와 문서 해시를 비공개 증인으로 사용하더라도, 외부 DPP가
논스를 공개한 이후의 비밀성을 주장하지 않는다. 후보 문서 값
이 알려진 경우 공개 논스가 후보 대입을 막아 주지도 않는다.
검증 범위는 세 문서 필드와 정책의 결합으로 제한되며, 그 밖의
DPP 내용이나 실제 제품의 안전성까지 인증하지 않는다.

% Source PDF page 9, column 2
